\documentclass[11pt,a4paper]{article}
\usepackage[utf8]{inputenc}
\usepackage[T1]{fontenc}
\usepackage[brazil]{babel}
\usepackage{geometry}
\usepackage{listings}
\usepackage{hyperref}
\geometry{margin=2.3cm}
\lstset{basicstyle=\ttfamily\small,breaklines=true,columns=fullflexible,frame=single,showstringspaces=false}
\title{Soluções dos exercícios\\Data Analysis with Python 2024--2025 -- Parte 1}
\author{Soluções independentes baseadas nos enunciados do PDF fornecido}
\date{2026}
\begin{document}
\maketitle
\tableofcontents
\section{Exercício 1.1 -- hello world}
\begin{lstlisting}[language=Python]
print("Hello, world!")
\end{lstlisting}

\section{Exercício 1.2 -- compliment}
\begin{lstlisting}[language=Python]
country = input("What country are you from? ")
print(f"I have heard that {country} is a beautiful country.")
\end{lstlisting}

\section{Exercício 1.3 -- multiplication}
\begin{lstlisting}[language=Python]
for i in range(11):
    print(f"4 multiplied by {i} is {4*i}")
\end{lstlisting}

\section{Exercício 1.4 -- multiplication table}
\begin{lstlisting}[language=Python]
for i in range(1, 11):
    for j in range(1, 11):
        print(f"{i*j:4}", end="")
    print()
\end{lstlisting}

\section{Exercício 1.5 -- two dice}
\begin{lstlisting}[language=Python]
for a in range(1, 7):
    for b in range(1, 7):
        if a + b == 5:
            print((a, b))
\end{lstlisting}

\section{Exercício 1.6 -- triple square}
\begin{lstlisting}[language=Python]
def triple(x):
    return 3*x

def square(x):
    return x*x

for i in range(1, 11):
    if square(i) > triple(i):
        break
    print(f"triple({i})=={triple(i)} square({i})=={square(i)}")
\end{lstlisting}

\section{Exercício 1.7 -- areas of shapes}
\begin{lstlisting}[language=Python]
import math

while True:
    shape = input("Choose a shape (triangle, rectangle, circle): ")
    if shape == "":
        break

    if shape == "triangle":
        base = float(input("Give base of the triangle: "))
        height = float(input("Give height of the triangle: "))
        area = base * height / 2
    elif shape == "rectangle":
        width = float(input("Give width of the rectangle: "))
        height = float(input("Give height of the rectangle: "))
        area = width * height
    elif shape == "circle":
        radius = float(input("Give radius of the circle: "))
        area = math.pi * radius * radius
    else:
        print("Unknown shape!")
        continue

    print(f"The area is {area:f}")
\end{lstlisting}

\section{Exercício 1.8 -- solve quadratic}
\begin{lstlisting}[language=Python]
from math import sqrt

def solve_quadratic(a, b, c):
    d = b*b - 4*a*c
    return ((-b + sqrt(d)) / (2*a), (-b - sqrt(d)) / (2*a))
\end{lstlisting}

\section{Exercício 1.9 -- merge}
\begin{lstlisting}[language=Python]
def merge(L1, L2):
    result = []
    i = j = 0
    while i < len(L1) and j < len(L2):
        if L1[i] <= L2[j]:
            result.append(L1[i])
            i += 1
        else:
            result.append(L2[j])
            j += 1
    result.extend(L1[i:])
    result.extend(L2[j:])
    return result
\end{lstlisting}

\section{Exercício 1.10 -- detect ranges}
\begin{lstlisting}[language=Python]
def detect_ranges(L):
    values = sorted(L)
    result = []
    i = 0
    while i < len(values):
        start = values[i]
        end = start
        while i + 1 < len(values) and values[i+1] == values[i] + 1:
            i += 1
            end = values[i]
        if end > start:
            result.append((start, end + 1))
        else:
            result.append(start)
        i += 1
    return result
\end{lstlisting}

\section{Exercício 1.11 -- interleave}
\begin{lstlisting}[language=Python]
def interleave(*lists):
    result = []
    for group in zip(*lists):
        result.extend(group)
    return result
\end{lstlisting}

\section{Exercício 1.12 -- distinct characters}
\begin{lstlisting}[language=Python]
def distinct_characters(L):
    return {s: len(set(s)) for s in L}
\end{lstlisting}

\section{Exercício 1.13 -- reverse dictionary}
\begin{lstlisting}[language=Python]
def reverse_dictionary(d):
    result = {}
    for english, finnish_words in d.items():
        for finnish in finnish_words:
            result.setdefault(finnish, []).append(english)
    return result
\end{lstlisting}

\section{Exercício 1.14 -- find matching}
\begin{lstlisting}[language=Python]
def find_matching(L, search):
    return [i for i, value in enumerate(L) if search in value]
\end{lstlisting}

\section{Exercício 1.15 -- two dice comprehension}
\begin{lstlisting}[language=Python]
pairs = [(a, b) for a in range(1, 7) for b in range(1, 7) if a + b == 5]
for pair in pairs:
    print(pair)
\end{lstlisting}

\section{Exercício 1.16 -- transform}
\begin{lstlisting}[language=Python]
def transform(s1, s2):
    a = map(int, s1.split())
    b = map(int, s2.split())
    return [x*y for x, y in zip(a, b)]
\end{lstlisting}

\section{Exercício 1.17 -- positive list}
\begin{lstlisting}[language=Python]
def positive_list(L):
    return list(filter(lambda x: x > 0, L))
\end{lstlisting}

\section{Exercício 1.18 -- acronyms}
\begin{lstlisting}[language=Python]
import string

def acronyms(text):
    result = []
    for word in text.split():
        cleaned = word.strip(string.punctuation)
        if len(cleaned) >= 2 and cleaned.isupper():
            result.append(cleaned)
    return result
\end{lstlisting}

\section{Exercício 1.19 -- sum equation}
\begin{lstlisting}[language=Python]
def sum_equation(L):
    if not L:
        return "0 = 0"
    return " + ".join(map(str, L)) + f" = {sum(L)}"
\end{lstlisting}

\section{Exercício 1.20 -- usemodule}
\subsection*{triangle.py}
\begin{lstlisting}[language=Python]
"""Funcoes para triangulos retangulos."""
from math import hypot

__version__ = "1.0"
__author__ = "Autor"

def hypotenuse(a, b):
    """Retorna o comprimento da hipotenusa."""
    return hypot(a, b)

def area(a, b):
    """Retorna a area de um triangulo retangulo."""
    return a * b / 2
\end{lstlisting}

\subsection*{usemodule.py}
\begin{lstlisting}[language=Python]
import triangle

def main():
    print(triangle.hypotenuse(3, 4))
    print(triangle.area(3, 4))

if __name__ == "__main__":
    main()
\end{lstlisting}

\section{Observação sobre as soluções}
As soluções acima são implementações independentes construídas a partir dos enunciados presentes no PDF fornecido. Elas não são uma cópia de um gabarito oficial do MOOC.
\end{document}
